% 来源：OEBPS/Text/01_14.xhtml（OPF spine 顺序；仅第一册）
\cleardoublepage
\BookEntry[牛顿《自然哲学的数学原理》]{牛顿《自然哲学的数学原理》}

\BookCover{images/书单01自然哲学的数学原理.jpg}{[英] 牛顿著| 赵振江译| 商务印书馆2006年版}

\begin{BookSummaryBox}
\textbf{标志着经典力学体系的建立，\\是人类科学史乃至整个人类文明史中的不朽巨著}
\par
\end{BookSummaryBox}

\subsubsection*{这本书谈什么？}

在牛顿之前，有伽利略、笛卡尔、开普勒三位科学先驱。伽利略提出了运动的相对性、运动的合成和自由落体的法则；笛卡尔讨论了惯性原理的动量守恒法则；开普勒提出了行星运动三定律。
\par

不过在牛顿之前，天体运动和地上运动是被分开来考虑的。牛顿的主要工作是将天体和地上的运动研究成果统一起来，并从数学上给予适当的证明。《自然哲学的数学原理》给出了著名的牛顿运动三大定律以及引力和距离的平方反比关系，构建了一个完整的经典力学体系。全书按照几何学的形式，以一种公理化的体系写成。科学史和人类文明史上，《自然哲学的数学原理》占据了极其重要的位置。毫不夸张地说，没有《自然哲学的数学原理》，就没有现代工业文明。
\par

\subsubsection*{芒格为什么谈起这本书？}

1986年6月13日，芒格在哈佛学校的毕业演讲中，提到了会使得生活痛苦的四个“忠告”。芒格的第二个“忠告”是：“尽可能从你们自身的经验获得知识，尽量别从其他人成功或失败的经验中广泛地吸取教训，不管他们是古人还是今人。”芒格谈道：
\par

\begin{quotation}
如果我再讲一个简短的历史故事，或许你们可以看得更清楚，从而更有效地过上与幸福无缘的生活。从前有个人，他勤奋地掌握了前人最优秀的成果，尽管开始研究分析几何的时候他的基础并不好，学得非常吃力。最终，他本人取得的成就引起了众人的瞩目，他是这样评价他自己的成果的：
\par

“如果说我比其他人看得更远，那是因为我站在巨人的肩膀上。”
\par

这人的骨灰如今埋在威斯敏斯特大教堂里，他的墓碑上有句异乎寻常的墓志铭：
\par

“这里安葬着永垂不朽的艾萨克·牛顿爵士。”
\par
\end{quotation}

\subsubsection*{《自然哲学的数学原理》作者牛顿如是说}

“圆滑是一门表达观点而又不树敌的艺术。”
\par

“如果说我曾发现什么有价值的原理，那应该归功于我的耐心观察，而非其他才能。”
\par

“没有大胆的猜想，就没有伟大的发现。”
\par

“我们造的墙太多，而修的桥太少。”
\par

“如果我算是一个人物（这一点我非常怀疑）的话，那么我的地位是通过努力工作而获得的。”
\par

“我总是把我的研究主题摆在面前，耐心等候，直到第一缕晨光初现，并逐渐变成一片光明。”
\par

“如果说我对公众有什么贡献的话，那是因为我的耐心思考。”
\par

“如果说我比其他人看得更远，那是因为我站在巨人的肩膀上。”
\par
